%----------------------------------------------------------------------------
\appendix
%----------------------------------------------------------------------------
\chapter*{\fuggelek}\addcontentsline{toc}{chapter}{\fuggelek}
\setcounter{chapter}{\appendixnumber}
%\setcounter{equation}{0} % a fofejezet-szamlalo az angol ABC 6. betuje (F) lesz
\numberwithin{equation}{section}
\numberwithin{figure}{section}
\numberwithin{lstlisting}{section}
%\numberwithin{tabular}{section}

\section{A Tesztadat-generátor Forráskódja (Python)}
\label{app:python_generator}

Az alábbi Python szkript felelős a szintetikus tesztszámlák (PDF) előállításáért. A kód az \texttt{FPDF} könyvtárat használja, és három különböző sablon (\textit{Modern, Simple, General}) alapján képes dinamikusan generálni a dokumentumokat.

\begin{lstlisting}[language=Python, caption={A ComplexInvoice osztály és a PDF generáló logika}, label={code:app_python_full}]
from fpdf import FPDF
import random

class ComplexInvoice(FPDF):
    def header(self):
        # Fejlec generalasa (Logo, Cegadatok)
        self.set_font('Arial', 'B', 12)
        self.cell(0, 10, 'INVOICE', 0, 1, 'C')
        self.ln(10)

    def footer(self):
        # Lablec (Oldalszam)
        self.set_y(-15)
        self.set_font('Arial', 'I', 8)
        self.cell(0, 10, f'Page {self.page_no()}', 0, 0, 'C')

    def generate_invoice(self, data, template_type="General"):
        self.add_page()
        self.set_font('Arial', '', 12)
        
        # Szallito es Vevo adatok
        self.cell(0, 10, f"Sender: {data.get('sender_text')}", 0, 1)
        self.cell(0, 10, f"Bill To: {data.get('bill_to_text')}", 0, 1)
        self.ln(5)
        
        # Szamla reszletei (Tablazatos forma)
        self.set_fill_color(200, 220, 255)
        self.cell(40, 10, 'Description', 1, 0, 'C', 1)
        self.cell(30, 10, 'Quantity', 1, 0, 'C', 1)
        self.cell(30, 10, 'Unit Price', 1, 0, 'C', 1)
        self.cell(30, 10, 'Total', 1, 1, 'C', 1)
        
        for item in data.get('items', []):
            self.cell(40, 10, str(item['desc']), 1)
            self.cell(30, 10, str(item['qty']), 1)
            self.cell(30, 10, str(item['unit']), 1)
            self.cell(30, 10, str(item['total']), 1, 1)
            
        self.ln(10)
        self.set_font('Arial', 'B', 12)
        self.cell(0, 10, f"Total Amount: {data.get('total_gross')} {data.get('currency')}", 0, 1, 'R')

# Pelda hivas:
# pdf = ComplexInvoice()
# pdf.generate_invoice(base_data)
# pdf.output("test_invoice.pdf")
\end{lstlisting}

\section{SAP Build Process Automation Szkriptek}
\label{app:bpa_scripts}

A prototípusban használt egyedi JavaScript kódok, amelyek a folyamatvezérlés (Automation) során futnak le.

\subsection{Adatfeltöltés és API Hívás (Prepare Upload)}

Ez a függvény készíti elő a \texttt{multipart/form-data} kérést a DOX szolgáltatás felé.

\begin{lstlisting}[language=JavaScript, caption={A dokumentumfeltöltést előkészítő JavaScript kód}, label={code:app_js_upload}]
// Input: filePath, accessToken
var myToken = "Bearer " + accessToken;

var optionsPayload = {
    documentType: "custom",
    schemaName: "SAP_invoice_schema",
    clientId: "default"
};

return {
    method: "POST",
    url: "https://aiservices-dox.cfapps.eu10.hana.ondemand.com/document-information-extraction/v1/document/jobs",
    responseType: "json",
    resolveBodyOnly: true,
    metadataType: irpa_core.enums.request.metadataType.formData, 
    headers: {
        Authorization: myToken
    },
    metadata: [
        { name: "file", file: filePath },
        { name: "options", value: JSON.stringify(optionsPayload) }
    ]
};
\end{lstlisting}

\subsection{Adatkinyerés és Normalizálás (Extract Data)}

A DOX szolgáltatásból érkező nyers JSON válasz feldolgozása és egyszerűsített objektummá alakítása.

\begin{lstlisting}[language=JavaScript, caption={A JSON válasz feldolgozása és normalizálása}, label={code:app_js_extract}]
// Input: apiResult
var extracted = {
    VendorName: "Unknown",
    InvoiceNumber: "",
    GrossAmount: 0,
    Currency: "EUR"
};

if (apiResult && apiResult.extraction) {
    var fields = apiResult.extraction.headerFields;
    for (var i = 0; i < fields.length; i++) {
        var name = fields[i].name;
        var value = fields[i].value;
        
        if (name === "senderName") extracted.VendorName = value;
        if (name === "documentNumber") extracted.InvoiceNumber = value;
        if (name === "totalAmount") extracted.GrossAmount = value;
        if (name === "currencyCode") extracted.Currency = value;
    }
}
return extracted;
\end{lstlisting}

\section{Tesztelési Jegyzőkönyv}
\label{app:test_logs}

Az 5. fejezetben bemutatott tesztesetek részletes futási naplója.

\begin{table}[H]
\centering
\begin{tabular}{|l|l|l|l|}
\hline
\textbf{Dátum} & \textbf{Teszteset} & \textbf{Eredmény} & \textbf{Megjegyzés} \\ \hline
2025.10.20 & TC\_01 (Touchless) & SIKERES & Automatikus jovahagyas megtortent (45s) \\ \hline
2025.10.20 & TC\_02 (Missing ID) & SIKERES & Validacios hiba, elutasito urlap generalva \\ \hline
2025.10.21 & TC\_03 (High Value) & SIKERES & Jovahagyasi feladat megjelent a Work Zone-ban \\ \hline
2025.10.21 & TC\_03 (Approval) & SIKERES & Felhasznaloi jovahagyas utan parkolas OK \\ \hline
\end{tabular}
\caption{Összesített tesztelési napló (2025. október)}
\label{tab:app_test_log}
\end{table}